%==============================================================================
% PARTIDA 08: SISTEMA DE PUESTA A TIERRA
% Codigo: OE.5.4.1.1
%==============================================================================
\subsection{OE.5.4.1.1 Puesta a Tierra PAT-1}

\subsubsection{Especificaciones Tecnicas}

El sistema de puesta a tierra (PAT-1) comprende la instalacion del electrodo vertical, pozo de tierra, conductor de proteccion, caja de registro y accesorios necesarios para proporcionar una via de baja impedancia a las corrientes de falla. Debe cumplir con lo establecido en el CNE-U Seccion 6 y la Norma EM.010.

El sistema consiste en una varilla de cobre (copperweld) de 5/8'' de diametro x 2.40 m de longitud, instalada verticalmente en un pozo de tierra. La resistencia de puesta a tierra debe ser menor a 15 Ohms, segun recomendacion para instalaciones residenciales.

\subsubsection{Requisitos de Materiales}

\begin{itemize}
  \item Varilla de cobre (copperweld) de 5/8'' x 2.40 m.
  \item Conector split-bolt de bronce para conexion varilla-conductor.
  \item Conductor de cobre temple blando 6 mm2 (desnudo) para conexion desde varilla a barra de tierra.
  \item Caja de registro de concreto de 0.30 m x 0.30 m x 0.30 m, con tapa de concreto.
  \item Gel conductor para mejorar el contacto varilla-tierra.
  \item Bentonita o tierra fertil para relleno del pozo.
  \item Tuberia PVC de 3/4'' para proteccion del conductor en tramos expuestos.
  \item Terminal de compression tipo ojo para conexion a barra de tierra.
\end{itemize}

\subsubsection{Metodo de Ejecucion}

\begin{enumerate}
  \item Ubicar el punto de instalacion del pozo de tierra segun plano IE-06, en zona accesible para inspeccion y mantenimiento, cercano al tablero general TG-01.
  \item Excavar un pozo de 0.50 m x 0.50 m x 2.50 m de profundidad.
  \item Colocar la varilla de cobre en el centro del pozo y hincarla verticalmente hasta una profundidad de 2.40 m, dejando 0.10 m sobresaliendo del fondo del pozo para conexiones.
  \item Conectar el conductor de cobre de 6 mm2 a la varilla mediante conector split-bolt de bronce, aplicando gel conductor en la conexion.
  \item Rellenar el pozo con tierra cernida mezclada con bentonita en capas de 0.20 m, compactando cada capa.
  \item Instalar la caja de registro de concreto sobre el pozo, de modo que permita el acceso a la conexion varilla-conductor.
  \item Tender el conductor de 6 mm2 desde la caja de registro hasta la barra de tierra del TG-01, protegido con tuberia PVC en tramos expuestos.
  \item Conectar el conductor a la barra de tierra del tablero mediante terminal de compression tipo ojo.
  \item Sellar la conexion con gel dielectrico para evitar corrosion.
\end{enumerate}

\subsubsection{Control de Calidad}

\begin{itemize}
  \item Medir la resistencia de puesta a tierra con telurómetro (debe ser menor a 15 Ohms).
  \item Verificar la continuidad electrica entre la barra de tierra del tablero y la varilla.
  \item Inspeccionar la conexion split-bolt (debe estar firme y con gel conductor).
  \item Comprobar que no existan empalmes en el conductor de tierra.
  \item Verificar que la caja de registro quede accesible y con tapa.
  \item Realizar medicion de resistencia en epoca seca y humeda si es posible.
\end{itemize}

\subsubsection{Metodo de Medicion}

Se medira por juego (jgo) de puesta a tierra completamente instalado, incluyendo varilla, conductor, caja de registro, conectores y accesorios, con medicion de resistencia verificada.

\subsubsection{Unidad de Medida}

jgo (juego) -- Incluye todos los componentes del sistema de puesta a tierra.

\subsubsection{Forma de Pago}

Pago por juego completamente instalado, con medicion de resistencia menor a 15 Ohms verificada por supervision, incluyendo materiales, mano de obra y pruebas.

%==============================================================================
% PARTIDA 09: ARTEFACTOS DE ALUMBRADO
% Codigo: OE.5.5.1
%==============================================================================
\subsection{OE.5.5.1 Luminaria LED Interior de Adosar 12 W}

\subsubsection{Especificaciones Tecnicas}

Luminaria LED de 12 W para iluminacion general interior, tipo adosar a techo, con temperatura de color 3000 K a 4000 K (blanco calido a neutro), eficiencia luminosa minima de 100 lm/W, y factor de potencia mayor a 0.90. Debe incluir driver LED integrado y difusor de policarbonato opal.

La luminaria debe contar con certificacion de producto y cumplir con las normas NTP correspondientes.

\subsubsection{Requisitos de Materiales}

\begin{itemize}
  \item Luminaria LED de adosar 12 W, incluye driver LED.
  \item Difusor de policarbonato opal.
  \item Tornillos de fijacion y tarugos de expansion para montaje en techo.
  \item Conectores rapidos tipo "wago" para conexion electrica.
  \item Cinta aislante.
\end{itemize}

\subsubsection{Metodo de Ejecucion}

\begin{enumerate}
  \item Verificar que la caja octogonal de salida de alumbrado este correctamente instalada y accesible.
  \item Retirar la tapa ciega de la caja octogonal.
  \item Identificar los conductores de fase, neutro y tierra en la caja de salida.
  \item Fijar la luminaria LED a la caja octogonal mediante tornillos.
  \item Conectar los conductores: fase (L) a fase, neutro (N) a neutro, y tierra (PE) si la luminaria lo requiere.
  \item Colocar el difusor y verificar el correcto sellado.
  \item Energizar el circuito y probar el funcionamiento de la luminaria.
\end{enumerate}

\subsubsection{Control de Calidad}

\begin{itemize}
  \item Verificar que la luminaria encienda correctamente sin parpadeos.
  \item Medir la tension de alimentacion en los bornes de la luminaria (220 V ± 10\%).
  \item Inspeccionar la correcta fijacion de la luminaria al techo.
  \item Verificar que el difusor este correctamente colocado y sellado.
\end{itemize}

\subsubsection{Metodo de Medicion}

Se medira por pieza (pza) de luminaria LED instalada y probada.

\subsubsection{Unidad de Medida}

pza (pieza).

\subsubsection{Forma de Pago}

Pago por pieza instalada, probada (encendido correcto) y aprobada por supervision.

%==============================================================================
% PARTIDA 10: EQUIPOS ELECTRICOS Y MECANICOS
% Codigo: OE.5.6.1
%==============================================================================
\subsection{OE.5.6.1 Electrobomba Monofasica de Agua 1 HP, 220V, 60Hz}

\subsubsection{Especificaciones Tecnicas}

Electrobomba monofasica centrifuga de 1 HP (745 W), para suministro de agua en la vivienda, con tension de alimentacion 220 V - 60 Hz. Debe incluir interruptor de nivel automatico (control de nivel) o presostato, segun el sistema de almacenamiento de agua de la vivienda.

La electrobomba debe contar con proteccion termica incorporada contra sobrecarga y cumplir con las normas NTP correspondientes para equipos electricos.

\subsubsection{Requisitos de Materiales}

\begin{itemize}
  \item Electrobomba monofasica centrifuga 1 HP, 220 V, 60 Hz.
  \item Interruptor de nivel automatico (control de nivel) o presostato.
  \item Cable de alimentacion con conector tipo industrial.
  \item Conductor de cobre LSOH 2.5 mm2 para alimentacion dedicada desde el tablero (OE.5.2.3.5).
  \item Interruptor termomagnetico 2P-16A para proteccion del circuito.
  \item Tomacorriente de servicio pesado 20 A (OE.5.2.1.2.3) para conexion del equipo.
  \item Válvula de compuerta y valvula de retencion para instalacion hidraulica.
  \item Accesorios de conexion hidraulica (nipples, uniones, cinturon teflon).
\end{itemize}

\subsubsection{Metodo de Ejecucion}

\begin{enumerate}
  \item Ubicar la electrobomba en el area tecnica designada (tercer piso o nivel de servicio), segun plano.
  \item Fijar la electrobomba a una base de concreto o soporte metalico mediante pernos de anclaje.
  \item Conectar la tuberia de succion y descarga de la electrobomba, utilizando cinturon teflon en las roscas.
  \item Instalar la valvula de compuerta en la salida y la valvula de retencion para evitar el retroceso del agua.
  \item Conectar el cable de alimentacion de la electrobomba al tomacorriente de servicio pesado 20 A.
  \item Conectar el control de nivel automatico segun las instrucciones del fabricante.
  \item Verificar que la proteccion termomagnetica del circuito sea de 16 A.
\end{enumerate}

\subsubsection{Control de Calidad}

\begin{itemize}
  \item Verificar que la electrobomba funcione correctamente (arranque y parada).
  \item Medir la corriente de operacion (debe ser menor a la corriente nominal de placa).
  \item Comprobar la correcta conexion del control de nivel automatico.
  \item Verificar que no existan fugas de agua en las conexiones hidraulicas.
  \item Medir la tension de alimentacion en bornes del equipo (220 V ± 10\%).
\end{itemize}

\subsubsection{Metodo de Medicion}

Se medira por unidad (UND) de electrobomba completamente instalada, incluyendo conexiones electricas, hidraulicas y pruebas de funcionamiento.

\subsubsection{Unidad de Medida}

UND (unidad).

\subsubsection{Forma de Pago}

Pago por unidad instalada, probada en operacion (arranque/parada) y aprobada por supervision.

%==============================================================================
% PARTIDA 11: PRUEBAS ELECTRICAS
% Codigo: OE.5.7.1
%==============================================================================
\subsection{OE.5.7.1 Pruebas de Aislamiento, Continuidad de Linea y Resistencia de Puesta a Tierra}

\subsubsection{Especificaciones Tecnicas}

Conjunto de pruebas electricas para verificar la correcta instalacion y seguridad de la instalacion electrica antes de la puesta en servicio. Incluye:

\begin{itemize}
  \item Prueba de continuidad de conductores de proteccion (tierra).
  \item Prueba de resistencia de aislamiento de conductores.
  \item Prueba de polaridad de tomacorrientes.
  \item Medicion de resistencia de puesta a tierra.
  \item Prueba de funcionamiento de interruptores diferenciales.
  \item Prueba de funcionamiento de interruptores termomagneticos.
  \item Prueba de tension de alimentacion en tableros y tomacorrientes.
\end{itemize}

Todas las pruebas deben realizarse segun la NTP 370.304 y la NTP-IEC 60364-6 (Verificacion de instalaciones electricas).

\subsubsection{Requisitos de Materiales y Equipos}

\begin{itemize}
  \item Megohmetro (Megger) de 500 V y 1000 V para prueba de aislamiento.
  \item Telurómetro para medicion de resistencia de puesta a tierra.
  \item Multimetro digital true RMS.
  \item Comprobador de tomacorrientes (tester de polaridad).
  \item Comprobador de interruptores diferenciales (RCD tester).
  \item Pinza amperimetrica para medicion de corriente.
  \item Protocolo de pruebas para registro de resultados.
\end{itemize}

\subsubsection{Metodo de Ejecucion}

\begin{enumerate}
  \item Antes de energizar, realizar prueba de continuidad de todos los conductores de proteccion a tierra.
  \item Medir resistencia de aislamiento entre fase-neutro, fase-tierra y neutro-tierra con megohmetro de 1000 V.
  \item Energizar la instalacion y medir tension en el tablero general (debe ser 220 V ± 10\%).
  \item Verificar la polaridad en cada tomacorriente con el comprobador de tomacorrientes.
  \item Probar cada interruptor diferencial presionando el boton TEST (debe disparar en menos de 40 ms).
  \item Medir la resistencia de puesta a tierra con telurómetro (metodo de 3 varillas auxiliares).
  \item Probar el funcionamiento de cada interruptor termomagnetico simulando sobrecarga (si es posible) o verificando su accion manual.
  \item Probar cada luminaria y cada interruptor de luz.
  \item Registrar todos los resultados en el protocolo de pruebas.
\end{enumerate}

\subsubsection{Control de Calidad}

\begin{itemize}
  \item Verificar que la resistencia de aislamiento sea superior a 1 MOhm.
  \item Comprobar que la resistencia de puesta a tierra sea menor a 15 Ohms.
  \item Verificar que la tension en todos los tomacorrientes sea 220 V ± 10\%.
  \item Confirmar que todos los interruptores diferenciales disparen correctamente con el boton TEST.
  \item Verificar que no exista continuidad entre neutro y tierra en los tomacorrientes.
  \item Entregar el protocolo de pruebas firmado por el responsable.
\end{itemize}

\subsubsection{Metodo de Medicion}

Se medira como un global (glb) que comprende la ejecucion de todas las pruebas descritas, incluyendo la emision del protocolo de pruebas.

\subsubsection{Unidad de Medida}

glb (global).

\subsubsection{Forma de Pago}

Pago por global de pruebas ejecutadas, con protocolo de pruebas presentado y aprobado por supervision.
